\Needspace{14\baselineskip}
\section{Discussion}
\label{sec:discussion}

% PRISMA Item 17: Discussion of results in the context of other evidence
% PRISMA Item 18: Limitations of the review

\subsection{Main findings}
\label{sec:principais-achados}

The 28 mapped studies reveal a field concentrated around one dominant
architecture: experience banks that distil resolution trajectories into
reusable entries (11 of 24 architecture studies; 46\%,
\autoref{tab:architectures}).
The LLM itself acts as the memory controller in 54\% of the
architecture studies.
Evaluation is centred on SWE-bench Verified, which appears in 13 of the
24 studies with benchmark evaluation (54\%, \autoref{tab:benchmarks}).
Studies increasingly choose SWE-bench Verified, but they do not test
memory designs against one another.
The 28 included studies are dominated by preprints (82\%) and use at
least eight distinct LLMs as backbones.
No study in the corpus compares three or more independent persistence
architectures under the same benchmark, model backbone, scaffold,
compute budget, and cost-accounting protocol.
The observed convergence reflects independent author choices rather
than an empirical validation of experience banks over alternatives such
as hierarchical knowledge bases or repository-history memory.

Controlled evidence is mostly positive but not uniformly so:
18 of the 19 architecture studies with a controlled no-persistence
baseline report positive aggregate gains in their main comparison
(\autoref{tab:performance}), while
CTIM-Rover~\citep{lindenbauer2025ctimrover} reports negative results
across its CTIM configurations;
within the SWE-bench Verified subset ($n=12$, including CTIM-Rover)
the median gain is 6.5~pp, and across the 18 comparable comparisons
the median is 7.8~pp, with reported gains ranging from $-$2.0~pp
to +39.0~pp.
The data are consistent with diminishing returns as the system
approaches the benchmark ceiling, though the pattern is descriptive
rather than formally tested: on SWE-bench Verified, the two largest
gains come from the two weakest baselines
(\citet{li2026traversalaspolicy}, 34.6\%~$\to$~+39.0~pp;
\citet{zhang2025darwingodel}, 20.0\%~$\to$~+30.0~pp),
and studies with baselines above 65\% report gains between +2.2 and
+9.4~pp~\citep{chen2025sweexp,wu2025gcc,xia2025livesweagent,deng2026memcoder}.
The cost impact is bimodal: systems that substitute condensed memory
for redundant context reduce token consumption by margins ranging from
about $-32$\% in total cost for
EET~\citep{guo2026eet} and $-39$\% in tokens for
GBT~\citep{li2026traversalaspolicy} to 2.4$\times$ fewer tokens for
KernelBlaster~\citep{dong2026kernelblaster},
while systems that add context via retrieval raise cost by +5 to
+21\%~\citep{chen2025sweexp,wang2026memgovern,xia2025livesweagent}.
TALM~\citep{shen2025talm} reports roughly half the tokens of MapCoder
on ClassEval, but that is a full-system comparison;
the memory module itself adds only marginal token overhead.

Adverse outcomes associated with persistence appear in nine of the 28
included studies (32\%, \autoref{tab:degradation}).
Four include direct performance drops or benchmark-level regressions;
the remaining five capture overhead, maintenance burden, security, and
evaluation artefacts.
Retrieval selection and backend governance are the dominant mechanisms
in the direct regressions:
\citet{lindenbauer2025ctimrover} report that unfiltered CTIM injection
biases exploration and falls up to 11~pp below the AutoCodeRover
baseline in its weakest configuration,
\citet{chen2025sweexp} report that retrieving more than one experience
($k > 1$) degrades performance from 42.0\% to 39.0\%, and
\citet{deshpande2025memtrack} identify that commercial memory backends
(Mem0, Zep) do not improve aggregate correctness over the no-memory
baseline, keep tool-call redundancy at roughly 21\%, and perform worse
on follow-up-question conditions.
The poisoning vulnerability documented by
\citet{srivastava2025memorygraft}, in which 9.1\% of poisoned
experiences appear in 47.9\% of retrieved records under dual-channel
retrieval (lexical + vector) while downstream behavioural severity is
not measured, indicates that memory stores without provenance checks
constitute a retrieval-level attack surface.

\subsection{Implications for research}
\label{sec:implicacoes-pesquisa}

The central gap identified by this review is the absence of head-to-head
comparisons.
No study compares three or more independent persistence architectures
under the same benchmark, model backbone, scaffold, compute budget, and
cost-accounting protocol.
Each system reports gains against its own no-memory baseline, but
experimental conditions differ in model, framework, instance subset, and
number of attempts.
This fragmentation prevents any conclusion on which architecture is
superior, because the observed gains may reflect implementation
differences as much as differences in memory architecture.
A comparative study evaluating at least three representative
architectures under controlled conditions would fill this gap.

Cost underreporting likewise limits the practical utility of the
results.
Only seven studies report monetary cost per task with enough
supporting detail for study-level cost
interpretation~\citep{chen2025sweexp,guo2026eet,wang2026memgovern,wu2025gcc,xia2025livesweagent,zhang2025accelopt,zhou2025tomswe},
and only a subset of these provide a within-study memory-vs-no-memory
cost delta;
eight of the 24 architecture studies (33\%) report no cost data, and
12 of the 28 included studies (43\%) do not contribute cost data when
contextual sources are counted, including the Du survey for which cost
is not applicable.
Without standardised cost-per-task metrics (tokens consumed, monetary
cost, wall-clock time), practitioners cannot assess the trade-off of
adopting persistence.
Constructing a cross-architecture Pareto frontier between accuracy and
cost would make adoption decisions measurable rather than accuracy-only.
ToM-SWE reports a configuration-level frontier within one architecture,
but the corpus does not contain a frontier that compares multiple
persistence designs under the same benchmark, model backbone, scaffold,
compute budget, and cost-accounting protocol.

Evidence on the relationship between architectural complexity and gain
is partial.
The data indicate that simple systems, such as
AccelOpt~\citep{zhang2025accelopt} with an experience queue and
EET~\citep{guo2026eet} with heuristic early termination, obtain
relevant gains in their respective domains without resorting to
multi-tier architectures.
However, this comparison is between distinct studies: benchmarks and
models differ across systems.
Only a study controlling the complexity variable internally could
confirm whether architectural simplicity is indeed sufficient or
whether scenarios exist where sophistication is justified.

Three studies propose memory as a directly executable artefact, an
emerging trend with still-limited evidence:
behaviour trees~\citep{li2026traversalaspolicy},
scripts created at run time~\citep{xia2025livesweagent}, and
whole agents~\citep{zhang2025darwingodel}.
The first two report substantial gains (+39.0~pp and +3.2~pp,
respectively), but the certainty of this finding is low given the
small number of studies and the absence of independent replication.

\subsection{Implications for practice}
\label{sec:implicacoes-pratica}

For teams testing persistence for the first time, the corpus supports
using a simple experience-bank implementation with calibrated retrieval
as a baseline, rather than adopting a more complex memory architecture
immediately.
Experience banks are the most frequent architecture in this corpus
(11 of 24 architecture studies; 46\%) and have reference open-source
implementations.
In high-baseline SWE-bench Verified settings ($>65\%$), reported gains
are modest but positive across several architecture families:
\citet{chen2025sweexp} and \citet{deng2026memcoder} report
experience-bank gains of $+2.2$ and $+9.4$~pp,
\citet{wu2025gcc} reports a Git-style version-control memory gain of
$+6.2$~pp, and \citet{xia2025livesweagent} reports $+3.2$~pp for
task-local executable tool creation.
These results support testing simple memory or context-management
baselines first, but they do not establish architectural superiority:
no study in the corpus compares experience banks against
hierarchical/state knowledge bases, repository-history memory, or other
architectures under controlled conditions.

Two ablations make retrieval calibration hard to ignore.
\citet{chen2025sweexp} report that retrieving more than one experience
($k > 1$) degrades performance from 42.0\% to 39.0\%, and
\citet{shen2026structurally} report that category isolation outperforms
global retrieval by 2.3~pp.
Both results point in the same direction: selectivity in retrieval
influences the effective gain, and excess retrieved context can
introduce noise.

For already-deployed systems, experience-guided early termination
offers cost reduction with a small reported accuracy cost.
\citet{guo2026eet} report 31.8\% savings in total cost with a maximum
loss of 0.2~pp, by identifying termination opportunities in 11.3\% of
instances.
This approach is agent-agnostic and requires no changes to the
underlying memory architecture, though the evidence rests on a single
study.

Of the 24 architecture studies, 17 share a complete repository and
three offer partial access, totalling 83\% with some degree of
reproducibility.

\subsection{Limitations of this review}
\label{sec:limitacoes}

The limitations of this review are organised into four categories of
validity threat.
Internal validity covers the rigour of the selection and extraction
processes.
External validity concerns generalisation of the findings beyond the
corpus.
Construct validity addresses the operationalisation of the measured
concepts.
Conclusion validity covers the stability of cross-study inferences.

Screening was conducted by a single researcher, which is an
acknowledged threat in systematic reviews and precludes the computation
of formal agreement indices such as Cohen's $\kappa$.
This limitation was mitigated by eligibility criteria calibrated in a
pilot study and refined through a stress test against five comparable
systematic reviews in
LLM/SE~\citep{hou2024llmse,zhang2024unifying,jin2024llmagents,wang2024llmagents,liu2024llmagents},
a process that produced 17 documented borderline cases and a decision
flow chart with eight sequential questions.
Each screening decision was recorded in a versioned spreadsheet with
per-criterion rationale, permitting retrospective audit of the
exclusions.

Data extraction was performed by the same researcher, introducing risk
of interpretation bias in qualitative fields such as architecture-type
classification and the categorisation of adverse outcomes.
Four mitigations reduced this risk: (i) programmatic validation of 448
enum values (16 fields $\times$ 28 studies), yielding eight corrections
on 1{,}036 extracted values (0.77\%);
(ii) a spot-check of 40 interpretive fields across 10 selected studies,
as detailed in \autoref{sec:extraction};
(iii) per-study evidence summaries derived from the full text of each
included paper, used as a reference when drafting the manuscript's
interpretive claims;
and (iv) re-verification of the \textit{negative\_result} fields
against the original PDF.
Even with these mitigations, formal reproducibility of the
classification in interpretive fields was not assessed.

The preprint prevalence (82\% of the corpus) reduces confidence in this
review's findings.
The field of knowledge persistence for coding agents is recent, with
publications concentrated between 2024 and 2026, and only five studies
have been published in peer-reviewed venues: two at
conferences~\citep{wang2026improving,zhang2025darwingodel} and three at
workshops~\citep{deshpande2025memtrack,shen2025talm,lindenbauer2025ctimrover}.
Sensitivity analysis S1 (\autoref{sec:sensitivity}) shows that,
when preprints are excluded, no qualitative conclusion inverts, but SQ4
becomes preprint-sensitive:
detailed monetary cost evidence is largely removed, and the remaining
peer-reviewed cost evidence is partial.

Searches were limited to two databases (Scopus and arXiv), complemented
by citation chasing.
Studies published exclusively in uncovered bases (for example, IEEE
Xplore, ACM Digital Library) may have been missed.
The rationale for this choice, based on \citet{mourao2020hybrid}, is
documented in \autoref{sec:fontes}.

The 13 studies that use SWE-bench Verified differ in base model, agent
framework, instance subset, and run conditions (\autoref{sec:sq2}).
The construct ``agent performance'' is not operationalised uniformly
across studies, which limits direct comparability of resolution
percentages reported as if they were equivalent quantities.
The taxonomy proposed in this review for persistence architectures also
embeds interpretive choices: families such as ``hybrid'' group
heterogeneous mechanisms whose fine differentiation would require
metadata absent from published texts.
The inclusion of Live-SWE-Agent~\citep{xia2025livesweagent} as a
task-local executable-memory boundary case broadens the
representational-substrate taxonomy but should not be interpreted as
evidence for cross-session persistence.

Heterogeneity across studies precludes formal meta-analysis and limits
inferences to qualitative patterns synthesised narratively.
Classifying findings into evidence types A through C
(\autoref{tab:certainty}) makes the empirical support of each inference
explicit: type A for direct counts, type B for aggregated cross-study
patterns, and type C for mechanism characterisations.
The two sensitivity analyses conducted (S1, preprint exclusion;
S2, exclusion of studies scoring below 3.0) did not invert any
quantitative finding, reinforcing the stability of the synthesis under
alternative inclusion criteria.

\subsection{Research agenda}
\label{sec:agenda-pesquisa}

The next empirical step should compare representative memory designs
under the same conditions rather than add another isolated system
report.
The absence of head-to-head comparisons calls for an experiment that
compares three or more independent persistence architectures under the
same benchmark, model backbone, scaffold, compute budget, and
cost-accounting protocol.
Cost underreporting requires that such an experiment include an
integrated cost-per-task analysis, producing a cross-architecture
Pareto frontier between accuracy and cost.

Metric standardisation is a prerequisite for future comparability.
The continual-learning metrics proposed by
\citet{joshi2025swebenchcl} (\textit{CL-Score},
\textit{forward/backward transfer}, \textit{forgetting}) have been
adopted by no other study in the corpus.
Adopting standardised cost metrics (tokens per task, monetary cost,
wall-clock time) and knowledge-retention metrics would enable
cross-study comparisons that are currently unfeasible.

Cross-agent sharing remains underexplored.
\citet{chen2025sweexp} is the only study with controlled evidence of
memory sharing, and even there it is limited to experience reuse within
the same framework;
\citet{tablan2025smarter} proposes a collective shared-memory layer for
coding agents but evaluates only synthetic experience under LLM-as-judge
metrics, leaving controlled effectiveness unestablished.
Sharing protocols across heterogeneous agents remain without empirical
evaluation.
The vulnerabilities flagged by \citet{srivastava2025memorygraft}
indicate that any such protocol must address the risk of poisoning in
open stores before it can be recommended for production.
